\section{清初地方治理、身份与社区证据}
清初征服的军报、诏令和接收文书密集，却不能把城池易手、薙发、迁界、驻防和设治直接等同为社区生活已经被重组。本节将所有清初主事件置于户口、旗籍、宗族、海岸迁徙、宗教与地方中介的材料问题中；应当承认，普通家庭和妇女的自述留存远少于清廷、南明和士人文本。
\subsection{1640—1649：地方接收、军事移动与社会证据}
下列节点按账本顺序排列。每项政策、战役或制度变化，都应与其在家庭、劳动、道路和身份分类上的不均衡后果区分开来。
\hypertarget{social:EQ-1644-SHANHAIGUAN}{}\paragraph{顺治元年四月（1644年5月，换算待核）}清军与吴三桂在山海关击败李自成军，打开进入华北的军事通道。地方层面应追问此事如何影响户口、迁徙、身份、劳役、土地或社区安全，而不是把行政分类当作群体自身的说明。核心来源为《清世祖实录》顺治元年四月山海关军报；《清史稿·世祖本纪》；《清初内国史院满文档案》顺治元年四月关宁军务档。编年记录定位朝廷所知的时间，不足以呈现城市、乡村、妇女、儿童和流动人口的全部经验。账本的研究方向是关宁军、满洲八旗与大顺军的协作及乞师叙事文本层累。可核条件为李自成入北京后关宁军与大顺在兵饷、统属和互信上破裂，多尔衮把危机转为入关机会。应分层追问的后果为清军获得穿越关隘的行动空间，北京政权争夺转为华北与南方的多战区征服。证据界限：原有置信注记：核心战事可靠；吴三桂求援文书原貌、各军兵数和战场过程受清廷与后出叙事塑形
\hypertarget{social:EQ-1644-BEIJING-ENTRY}{}\paragraph{顺治元年五月（1644年6月前后）}多尔衮率军进入北京并接收明朝京城的行政与礼制空间。地方层面应追问此事如何影响户口、迁徙、身份、劳役、土地或社区安全，而不是把行政分类当作群体自身的说明。核心来源为《清世祖实录》顺治元年五月；《清史稿·世祖本纪》；《甲申传信录》及北京地方记事。编年记录定位朝廷所知的时间，不足以呈现城市、乡村、妇女、儿童和流动人口的全部经验。账本的研究方向是北京接收、明清鼎革中的城市秩序和遗臣选择。可核条件为大顺撤离与明廷崩溃留下军政真空，清军需获取京师粮饷、官僚和象征中心。应分层追问的后果为清廷以北京为治所，开始将军事占领转化为对明朝制度资源的继承和重组。证据界限：原有置信注记：入城与接收可靠；城内秩序、迎降者范围和掠夺状况须以不同见证互校
\hypertarget{social:EQ-1644-SHUNZHI-BEIJING}{}\paragraph{顺治元年九月（1644年10月前后）}顺治帝自盛京抵北京，清廷举行定都与皇帝入城仪式。地方层面应追问此事如何影响户口、迁徙、身份、劳役、土地或社区安全，而不是把行政分类当作群体自身的说明。核心来源为《清世祖实录》顺治元年九月；《清史稿·世祖本纪》；《清初内国史院满文档案》北京迁都礼仪档。编年记录定位朝廷所知的时间，不足以呈现城市、乡村、妇女、儿童和流动人口的全部经验。账本的研究方向是盛京—北京双中心、迁都礼仪与征服合法性。可核条件为摄政政权虽占北京，仍需皇帝亲临与礼制重置来建立持续中央名义。应分层追问的后果为北京成为皇帝驻跸与部院中心，盛京仍保有满洲政治和象征地位。证据界限：原有置信注记：日期与仪式主干可靠；官修礼仪记录不能等同各地承认
\hypertarget{social:EQ-1644-REGENT-ORDERS}{}\paragraph{顺治元年冬（1644年）}多尔衮以摄政王名义整顿京师官署、招降明官并发布安民与征发命令。地方层面应追问此事如何影响户口、迁徙、身份、劳役、土地或社区安全，而不是把行政分类当作群体自身的说明。核心来源为《清世祖实录》顺治元年十月至十二月；《清史稿·职官志》；中国第一历史档案馆藏内阁大库顺治元年题本〔京师官署交接〕。编年记录定位朝廷所知的时间，不足以呈现城市、乡村、妇女、儿童和流动人口的全部经验。账本的研究方向是征服政权的文书接管、降官网络与命令执行断点。可核条件为新政权需要恢复行政传递、维持京师粮饷并吸纳明代文官技术。应分层追问的后果为中央可用的部院和任官网络开始恢复，但忠诚、财力与军纪继续限制覆盖面。证据界限：原有置信注记：命令发布可靠；接受任官和地方执行不能从诏令直接推出
\hypertarget{social:EQ-1645-NANJING-TAKEN}{}\paragraph{顺治二年四月（1645年5月前后）}多铎军攻取南京，南明弘光朝的都城失守。地方层面应追问此事如何影响户口、迁徙、身份、劳役、土地或社区安全，而不是把行政分类当作群体自身的说明。核心来源为《清世祖实录》顺治二年四月；《清史稿·世祖本纪》；《南明史料》所收弘光朝奏疏与江宁地方记事。编年记录定位朝廷所知的时间，不足以呈现城市、乡村、妇女、儿童和流动人口的全部经验。账本的研究方向是南京失守、江南军政资源与南明合法性危机。可核条件为清军沿运河与淮扬推进，弘光朝党争、军饷困难和江防失守相互叠加。应分层追问的后果为南明失去首都和财政文化中心，江浙、闽粤的抵抗及另立政权随之分化。证据界限：原有置信注记：城陷日期可靠；守军数量、降附动机和战后暴力须并读南明与地方材料
\hypertarget{social:EQ-1645-HONGGUANG-CAPTURE}{}\paragraph{顺治二年五月（1645年）}弘光帝朱由崧在芜湖一带被俘并被押往北京。地方层面应追问此事如何影响户口、迁徙、身份、劳役、土地或社区安全，而不是把行政分类当作群体自身的说明。核心来源为《清世祖实录》顺治二年五月；《清史稿·弘光帝本纪》；《明季南略》及江南地方志相关条目。编年记录定位朝廷所知的时间，不足以呈现城市、乡村、妇女、儿童和流动人口的全部经验。账本的研究方向是弘光朝崩解、地方军阀与降附网络。可核条件为南京失守后弘光朝没有持续的军事护卫与统一指挥，逃亡路线被清军和地方武装切断。应分层追问的后果为以弘光名义维系的朝廷失能，隆武、永历政权在东南和西南重建名义中心。证据界限：原有置信注记：被俘及押解主干可靠；参与者责任和移交地点在纪传、地方记事间有异
\hypertarget{social:EQ-1645-HAIR-ORDER}{}\paragraph{顺治二年六月（1645年）}清廷重申薙发易服命令，在江南等地引发服从、逃亡与抵抗。地方层面应追问此事如何影响户口、迁徙、身份、劳役、土地或社区安全，而不是把行政分类当作群体自身的说明。核心来源为《清世祖实录》顺治二年六月；《清史稿·舆服志》；内阁大库顺治二年兵科题本与江南府县志薙发条。编年记录定位朝廷所知的时间，不足以呈现城市、乡村、妇女、儿童和流动人口的全部经验。账本的研究方向是征服礼制、身体规训与地方反应差异。可核条件为清廷以发式服制识别新旧政治归属，军事推进又使命令与治安相互捆绑。应分层追问的后果为服制成为忠降识别和集体动员的触发点，若干江南城镇的冲突加剧征服成本。证据界限：原有置信注记：命令文本可靠；各地执行速度和社会反应高度不一，不能以诏书概括全国
\hypertarget{social:EQ-1645-YANGZHOU}{}\paragraph{顺治二年四—五月（1645年）}扬州攻陷后发生大规模杀戮与掳掠，规模和持续时间存在重大争议。地方层面应追问此事如何影响户口、迁徙、身份、劳役、土地或社区安全，而不是把行政分类当作群体自身的说明。核心来源为《清世祖实录》顺治二年四月扬州军报；《清史稿·史可法传》；王秀楚《扬州十日记》与《扬州府志》。编年记录定位朝廷所知的时间，不足以呈现城市、乡村、妇女、儿童和流动人口的全部经验。账本的研究方向是征服暴力的叙事、死亡估计与城市记忆。可核条件为扬州守军扼守清军南下通道，攻城后的军纪、报复和掠夺处于战时权力失控与奖惩机制中。应分层追问的后果为扬州经验在江南传播，既促成恐惧性降附，也强化若干城镇抵抗及后世纪念政治。证据界限：原有置信注记：扬州失陷可靠；死亡数字和“十日”叙事须标为后出见证并与地方材料互证
\hypertarget{social:EQ-1645-JIANGYIN}{}\paragraph{顺治二年闰六月至八月（1645年）}江阴在薙发命令和清军围攻中失守，地方记忆长期以守城与屠城叙述该事件。地方层面应追问此事如何影响户口、迁徙、身份、劳役、土地或社区安全，而不是把行政分类当作群体自身的说明。核心来源为《清世祖实录》顺治二年江南军报；《清史稿·世祖本纪》；《江阴县志》及《江阴城守纪》。编年记录定位朝廷所知的时间，不足以呈现城市、乡村、妇女、儿童和流动人口的全部经验。账本的研究方向是江南城镇抵抗、地方文本与屠城数字核验。可核条件为服制命令与守备危机把行政服从转成军事对抗，清军需确保沿江交通和粮道。应分层追问的后果为清廷以军事压制恢复沿江节点；江阴成为地方忠义记忆，伤亡规模不能由单一叙事定量。证据界限：原有置信注记：城陷主干可靠；围城天数、伤亡和人物言行多依赖地方后出文本
\hypertarget{social:EQ-1645-ZHEJIANG-CONQUEST}{}\paragraph{顺治二年六月至十月（1645年）}清军进取杭州、绍兴等地，南明鲁王政权转入浙东及海上活动。地方层面应追问此事如何影响户口、迁徙、身份、劳役、土地或社区安全，而不是把行政分类当作群体自身的说明。核心来源为《清世祖实录》顺治二年六月至十月；《清史稿·鲁王传》；《浙江通志》与鲁王政权文书辑录。编年记录定位朝廷所知的时间，不足以呈现城市、乡村、妇女、儿童和流动人口的全部经验。账本的研究方向是浙东海上抵抗、地方士绅与区域控制碎片化。可核条件为南京失守后浙东军民与明宗室试图保持独立指挥，清军沿水陆节点分割其资源。应分层追问的后果为清廷取得浙江大部城镇，鲁王集团与郑氏网络相连，海疆战争成为长期问题。证据界限：原有置信注记：城市失守与鲁王转移可靠；各府实际控制和海岛联系须逐地核查
\hypertarget{social:EQ-1646-LONGWU-CAPTURE}{}\paragraph{顺治三年八月（1646年）}清军在闽北俘杀隆武帝朱聿键，隆武朝终结。地方层面应追问此事如何影响户口、迁徙、身份、劳役、土地或社区安全，而不是把行政分类当作群体自身的说明。核心来源为《清世祖实录》顺治三年八月；《清史稿·隆武帝本纪》；《福建通志》与黄道周文集相关记载。编年记录定位朝廷所知的时间，不足以呈现城市、乡村、妇女、儿童和流动人口的全部经验。账本的研究方向是福建征服、隆武朝廷与郑氏的早期关系。可核条件为清军由浙赣入闽，隆武政权军饷和将领协调不足且难以掌握山海通道。应分层追问的后果为闽中明廷名义被打断，郑氏网络和地方武装成为东南抗清的重要行动者。证据界限：原有置信注记：被俘死亡主干可靠；行止路线与地方助力记载不一
\hypertarget{social:EQ-1646-SHAOWU-GUANGZHOU}{}\paragraph{顺治三年十一月（1646年）}清军攻入广州，绍武帝绍武帝被杀，广州短暂的绍武朝结束。地方层面应追问此事如何影响户口、迁徙、身份、劳役、土地或社区安全，而不是把行政分类当作群体自身的说明。核心来源为《清世祖实录》顺治三年十一月；《清史稿·绍武帝本纪》；《广东通志》及南明文集。编年记录定位朝廷所知的时间，不足以呈现城市、乡村、妇女、儿童和流动人口的全部经验。账本的研究方向是两广明廷竞争、广州港口与地方军政。可核条件为隆武朝覆灭后明宗室在广东争夺名义中心，内部冲突削弱对清军进攻的协调。应分层追问的后果为永历政权成为西南主要明廷名义，广东沿海郑氏与地方势力继续影响战局。证据界限：原有置信注记：广州失守可靠；绍武与永历两朝冲突细节主要来自后出南明叙事
\hypertarget{social:EQ-1646-YONGLI-ZHAOQING}{}\paragraph{永历元年／顺治三年（1646年冬）}朱由榔在肇庆即位为永历帝，建立延续至西南的南明朝廷。地方层面应追问此事如何影响户口、迁徙、身份、劳役、土地或社区安全，而不是把行政分类当作群体自身的说明。核心来源为《清世祖实录》顺治三年两广军报；《清史稿·永历帝本纪》；《永历实录》与《肇庆府志》。编年记录定位朝廷所知的时间，不足以呈现城市、乡村、妇女、儿童和流动人口的全部经验。账本的研究方向是永历朝形成、两广地方网络与南明多中心政治。可核条件为绍武朝覆灭及明宗室分散，迫使支持者在两广重建朝廷；西南军政资源提供退守条件。应分层追问的后果为永历年号成为西南抗清行动的共同名义，清军须把两广、贵州、云南与缅甸路线连为战场。证据界限：原有置信注记：即位时间大体可靠；改元与礼仪细节应按永历文书和清廷军报分别处理
\hypertarget{social:EQ-1648-JIN-SHENGHUAN}{}\paragraph{顺治五年（1648年）}南昌守将金声桓反清并改奉南明，江西战局重新打开。地方层面应追问此事如何影响户口、迁徙、身份、劳役、土地或社区安全，而不是把行政分类当作群体自身的说明。核心来源为《清世祖实录》顺治五年江西军报；《清史稿·金声桓传》；《南昌府志》与南明永历朝文书。编年记录定位朝廷所知的时间，不足以呈现城市、乡村、妇女、儿童和流动人口的全部经验。账本的研究方向是降将再选择、江西财政与南明联络。可核条件为清初任用降将以快速覆盖江南，但军饷、地位和地方网络未稳定；南明仍可提供合法性资源。应分层追问的后果为清廷被迫向江西投入围城和调兵资源，显示占领区行政接收不等于长期服从。证据界限：原有置信注记：反叛及南昌控制可靠；部众规模和改奉动机多为交战各方归责语言
\hypertarget{social:EQ-1648-LI-CHENGDONG}{}\paragraph{顺治五年（1648年）}广东清将李成栋反清，广州一度转入永历朝控制。地方层面应追问此事如何影响户口、迁徙、身份、劳役、土地或社区安全，而不是把行政分类当作群体自身的说明。核心来源为《清世祖实录》顺治五年广东军报；《清史稿·李成栋传》；《广东通志》及永历朝诏敕辑录。编年记录定位朝廷所知的时间，不足以呈现城市、乡村、妇女、儿童和流动人口的全部经验。账本的研究方向是广东降将政治、港口资源与永历朝复兴。可核条件为清军依靠原明将领维系两广，军功分配与中央控制的矛盾给永历朝重新联络的空间。应分层追问的后果为两广战场扩大并与江西、闽海相连，清廷须重新部署南方军饷和将领体系。证据界限：原有置信注记：反叛与广州易手可靠；个人动机和地方支持范围有强烈阵营叙事差异
\hypertarget{social:EQ-1649-NANCHANG}{}\paragraph{顺治六年正月至三月（1649年）}清军围攻并攻取南昌，金声桓反叛失败。地方层面应追问此事如何影响户口、迁徙、身份、劳役、土地或社区安全，而不是把行政分类当作群体自身的说明。核心来源为《清世祖实录》顺治六年正月至三月；《清史稿·金声桓传》；内阁大库顺治六年江西军需题本与《南昌府志》。编年记录定位朝廷所知的时间，不足以呈现城市、乡村、妇女、儿童和流动人口的全部经验。账本的研究方向是围城战、江西粮饷与征服后的城市恢复。可核条件为金声桓据南昌切断清廷赣江通道，清军需以持续军需和外围据点完成围攻。应分层追问的后果为江西反叛中心被击破，清廷恢复通道但须面对战后税源、流民和秩序重建。证据界限：原有置信注记：围城结局可靠；攻城损失、城内处置和战后人口变化须以地方资料复核
\subsection{1650—1659：地方接收、军事移动与社会证据}
下列节点按账本顺序排列。每项政策、战役或制度变化，都应与其在家庭、劳动、道路和身份分类上的不均衡后果区分开来。
\hypertarget{social:EQ-1650-GUANGZHOU}{}\paragraph{顺治七年十一月（1650年）}清军重新攻取广州，战后杀戮与损毁规模成为广东地方记忆中的重大争议。地方层面应追问此事如何影响户口、迁徙、身份、劳役、土地或社区安全，而不是把行政分类当作群体自身的说明。核心来源为《清世祖实录》顺治七年十一月；《清史稿·世祖本纪》；《广东通志》、广州地方笔记与内阁大库广东军报。编年记录定位朝廷所知的时间，不足以呈现城市、乡村、妇女、儿童和流动人口的全部经验。账本的研究方向是广州征服暴力、海港城市与地方记忆。可核条件为李成栋反叛后广东成为永历朝资源节点，清军重夺广州兼具军事和海防目的。应分层追问的后果为永历朝在两广基础削弱，广州损失与人口流动影响地方财政、宗族和港口恢复。证据界限：原有置信注记：广州再陷可靠；死亡数字和“庚寅之劫”持续时间不能从清军战报单定
\hypertarget{social:EQ-1650-DORGON-DEATH}{}\paragraph{顺治七年十二月（1650年）}摄政王多尔衮在古北口外行猎途中去世，摄政政治结构随之改组。地方层面应追问此事如何影响户口、迁徙、身份、劳役、土地或社区安全，而不是把行政分类当作群体自身的说明。核心来源为《清世祖实录》顺治七年十二月；《清史稿·多尔衮传》；《清初内国史院满文档案》摄政王丧仪与遗命档。编年记录定位朝廷所知的时间，不足以呈现城市、乡村、妇女、儿童和流动人口的全部经验。账本的研究方向是摄政体制、宗室权力与死后政治。可核条件为多尔衮集中军政与皇室事务，围绕摄政权的宗室和旗人竞争在其死后失去调停中心。应分层追问的后果为顺治帝及近臣得以重组朝廷，摄政集团的名誉与权力被重新界定。证据界限：原有置信注记：死亡日期可靠；死因、身后清算与权力斗争不能只依后来的政治定性
\hypertarget{social:EQ-1651-SHUNZHI-PERSONAL-RULE}{}\paragraph{顺治八年（1651年）}顺治帝亲政，撤除摄政安排并重置皇帝周边的决策网络。地方层面应追问此事如何影响户口、迁徙、身份、劳役、土地或社区安全，而不是把行政分类当作群体自身的说明。核心来源为《清世祖实录》顺治八年正月；《清史稿·世祖本纪》；《顺治朝内阁档》亲政诏令与任官题本。编年记录定位朝廷所知的时间，不足以呈现城市、乡村、妇女、儿童和流动人口的全部经验。账本的研究方向是顺治亲政、皇帝近臣与满汉官僚协作。可核条件为多尔衮去世使摄政权威无法延续，皇帝及支持者以礼制和任免重建指挥链。应分层追问的后果为中央命令署名和人事结构改变，但南方征服与财政仍要求依靠地方将领和降官。证据界限：原有置信注记：亲政事实可靠；实际决策仍受宗室、旗务和满汉官僚共同制约
\hypertarget{social:EQ-1652-DINGGUO-GUILIN}{}\paragraph{永历六年／顺治九年（1652年）}李定国攻克桂林，孔有德自尽，清在广西的军事据点受重创。地方层面应追问此事如何影响户口、迁徙、身份、劳役、土地或社区安全，而不是把行政分类当作群体自身的说明。核心来源为《清世祖实录》顺治九年广西军报；《清史稿·孔有德传》；《永历实录》与《广西通志》。编年记录定位朝廷所知的时间，不足以呈现城市、乡村、妇女、儿童和流动人口的全部经验。账本的研究方向是李定国军事行动、广西城防与南明短暂反攻。可核条件为两广清军战线长、补给困难且将领协调不足，永历军利用机动兵力突击桂林。应分层追问的后果为永历朝获得声望与物资，但未能稳定保持两广；清廷随即加强西南调兵。证据界限：原有置信注记：桂林战果可靠；双方兵数和孔有德死因细节应保留不同文本说法
\hypertarget{social:EQ-1652-DINGGUO-HENGZHOU}{}\paragraph{永历六年／顺治九年夏（1652年）}李定国在衡州击败尼堪军，清军西南攻势暂受挫。地方层面应追问此事如何影响户口、迁徙、身份、劳役、土地或社区安全，而不是把行政分类当作群体自身的说明。核心来源为《清世祖实录》顺治九年湖南军报；《清史稿·尼堪传》；《永历实录》与《衡州府志》。编年记录定位朝廷所知的时间，不足以呈现城市、乡村、妇女、儿童和流动人口的全部经验。账本的研究方向是湘桂战场、八旗机动军与南明补给。可核条件为清军试图由湖南压迫两广和贵州，李定国依托本地军队与道路选择反击。应分层追问的后果为清廷认识到西南征服不能只靠快速突进，转向更大规模调兵、饷运和分区推进。证据界限：原有置信注记：战役结局大体可靠；地点、伤亡与尼堪死亡过程在中清叙事与南明记录间有异
\hypertarget{social:EQ-1653-ZHENG-XIAMEN}{}\paragraph{顺治十年（1653年）}郑成功以厦门、金门为基地整合海上军政与贸易网络，持续攻击福建沿海清军据点。地方层面应追问此事如何影响户口、迁徙、身份、劳役、土地或社区安全，而不是把行政分类当作群体自身的说明。核心来源为《清世祖实录》顺治十年福建海防军报；《清史稿·郑成功传》；《厦门志》与郑氏文书辑录。编年记录定位朝廷所知的时间，不足以呈现城市、乡村、妇女、儿童和流动人口的全部经验。账本的研究方向是郑氏海上政权、闽南港口与军商网络。可核条件为隆武朝覆灭后郑成功保留明室名义，并借岛屿港口、商船和侨民网络筹饷。应分层追问的后果为福建海疆无法由陆上州县单独控制，清廷日后将海禁、迁界和水师建设结合。证据界限：原有置信注记：基地经营可靠；郑军船只、兵员和贸易收入多来自敌对战报或后出郑氏叙事
\hypertarget{social:EQ-1654-ZHENG-ZHANGZHOU}{}\paragraph{顺治十一年（1654年）}郑成功军围攻漳州等闽南城镇，清军守城与援军维持陆海据点。地方层面应追问此事如何影响户口、迁徙、身份、劳役、土地或社区安全，而不是把行政分类当作群体自身的说明。核心来源为《清世祖实录》顺治十一年福建军报；《清史稿·郑成功传》；《漳州府志》与荷兰东印度公司台湾商馆日记。编年记录定位朝廷所知的时间，不足以呈现城市、乡村、妇女、儿童和流动人口的全部经验。账本的研究方向是闽南城镇、海上补给与郑氏军事能力。可核条件为郑氏需扩大粮饷和招募基础，漳州连接闽南腹地与港口；清军依托城防和陆路援兵。应分层追问的后果为郑氏未能迅速夺取大陆府城，海岛基地与海上贸易的重要性进一步上升。证据界限：原有置信注记：围攻与未克主干可靠；郑氏作战目标和双方损失须以中荷航海记录交叉
\hypertarget{social:EQ-1656-ZHENG-QUANZHOU}{}\paragraph{顺治十三年（1656年）}郑成功军在泉州、同安一线进攻并争夺闽南水陆通道。地方层面应追问此事如何影响户口、迁徙、身份、劳役、土地或社区安全，而不是把行政分类当作群体自身的说明。核心来源为《清世祖实录》顺治十三年福建军报；《清史稿·郑成功传》；《泉州府志》及郑氏文书辑录。编年记录定位朝廷所知的时间，不足以呈现城市、乡村、妇女、儿童和流动人口的全部经验。账本的研究方向是闽南水陆通道、地方团练与郑氏征粮。可核条件为郑氏既需守厦金又要取得腹地粮食，清军以府县城防切断其陆上补给。应分层追问的后果为拉锯战使闽南居民承受征粮、迁徙与治安压力，也推动清廷酝酿强制沿海政策。证据界限：原有置信注记：战事发生可靠；具体攻守日期与城寨归属须按清军报、郑氏文书和地方志逐项核对
\hypertarget{social:EQ-1657-YONGLI-GUIZHOU}{}\paragraph{永历十一年／顺治十四年（1657年）}永历朝在贵州、云南的孙可望与李定国集团矛盾公开化，孙可望转降清廷。地方层面应追问此事如何影响户口、迁徙、身份、劳役、土地或社区安全，而不是把行政分类当作群体自身的说明。核心来源为《清世祖实录》顺治十四年贵州军报；《清史稿·孙可望传》；《永历实录》与贵州地方志。编年记录定位朝廷所知的时间，不足以呈现城市、乡村、妇女、儿童和流动人口的全部经验。账本的研究方向是西南军阀联盟、南明朝廷与降附政治。可核条件为长期战争中的饷源、官爵和指挥权竞争削弱永历朝合作，清廷提供孙可望转换阵营的条件。应分层追问的后果为清军获得西南军情与向导资源，永历朝失去大批可动员兵力，云南退路日益孤立。证据界限：原有置信注记：孙可望降清可靠；冲突责任、军队规模与永历帝处境受各方政治文本影响
\hypertarget{social:EQ-1658-QING-YUNNAN}{}\paragraph{顺治十五年（1658年）}清军由贵州、四川等方向进逼云南，永历朝廷撤离昆明并向西南退却。地方层面应追问此事如何影响户口、迁徙、身份、劳役、土地或社区安全，而不是把行政分类当作群体自身的说明。核心来源为《清世祖实录》顺治十五年云南军报；《清史稿·洪承畴传》；《云南通志》与永历朝奏疏辑录。编年记录定位朝廷所知的时间，不足以呈现城市、乡村、妇女、儿童和流动人口的全部经验。账本的研究方向是入滇战争、山地补给与地方政治。可核条件为孙可望降清削弱南明防线，清军依托多路会攻和地方土司、降将信息推进。应分层追问的后果为永历政权转为流动朝廷，清军虽占昆明等城，仍须在边地和交通线上持续作战。证据界限：原有置信注记：清军入滇和永历撤离可靠；各路行军与地方归附应分区而非视作全省一次完成
\hypertarget{social:EQ-1659-NANJING-EXPEDITION}{}\paragraph{顺治十六年七月至八月（1659年）}郑成功北上围攻南京，未能突破城防后撤回海上基地。地方层面应追问此事如何影响户口、迁徙、身份、劳役、土地或社区安全，而不是把行政分类当作群体自身的说明。核心来源为《清世祖实录》顺治十六年七月至八月；《清史稿·郑成功传》；《延平王户官杨英从征实录》与荷兰东印度公司海上报告。编年记录定位朝廷所知的时间，不足以呈现城市、乡村、妇女、儿童和流动人口的全部经验。账本的研究方向是郑成功北伐、长江防务与军事补给。可核条件为郑氏希望借长江下游粮赋和明遗民网络扩大反清战场，清军则集中守备江防和城门。应分层追问的后果为北伐失败暴露郑军内河攻坚和长期补给限制，促使其把台湾作为新的战略选择。证据界限：原有置信注记：北伐及撤退可靠；作战计划、内应和伤亡数字在清廷、郑氏与荷兰记录间冲突
\hypertarget{social:EQ-1659-YONGLI-MYANMAR}{}\paragraph{永历十三年／顺治十六年末（1659年）}永历帝一行进入缅甸，西南战争扩展为跨境避难和外交问题。地方层面应追问此事如何影响户口、迁徙、身份、劳役、土地或社区安全，而不是把行政分类当作群体自身的说明。核心来源为《清世祖实录》顺治十六年云南军报；《清史稿·永历帝本纪》；《琉璃宫史》缅文传统及欧洲旅行记译注。编年记录定位朝廷所知的时间，不足以呈现城市、乡村、妇女、儿童和流动人口的全部经验。账本的研究方向是永历入缅、明清战争与中缅边境政治。可核条件为清军推进使永历朝失去云南腹地，缅甸王朝面临接纳流亡朝廷与避免清军压力的两难。应分层追问的后果为永历政权补给与外交选择受缅甸宫廷制约，清廷将追捕与边境交涉结合。证据界限：原有置信注记：入缅事实可靠；路线、接待条件与缅甸朝廷决策须并读汉文和缅文/欧洲记载
\subsection{1660—1669：地方接收、军事移动与社会证据}
下列节点按账本顺序排列。每项政策、战役或制度变化，都应与其在家庭、劳动、道路和身份分类上的不均衡后果区分开来。
\hypertarget{social:EQ-1660-QING-YUNNAN-CONSOLIDATION}{}\paragraph{顺治十七年（1660年）}清军继续经营云南城镇与驿路，永历朝残部向中缅边地退缩。地方层面应追问此事如何影响户口、迁徙、身份、劳役、土地或社区安全，而不是把行政分类当作群体自身的说明。核心来源为《清世祖实录》顺治十七年云南军报；《清史稿·吴三桂传》；内阁大库顺治十七年云南粮饷题本与《云南通志》。编年记录定位朝廷所知的时间，不足以呈现城市、乡村、妇女、儿童和流动人口的全部经验。账本的研究方向是云南征服后的军需、土司关系与驿路。可核条件为占领昆明后清军仍须供应远征军、维持道路并争取地方中介，永历残部利用边地机动。应分层追问的后果为云南成为吴三桂等将领驻军和财政调度基地，也为后来的藩镇结构埋下条件。证据界限：原有置信注记：清军控制主要城镇大体可靠；土司、矿区和山地实际控制不能以府城易手推定
\hypertarget{social:EQ-1661-KANGXI-ACCESSION}{}\paragraph{顺治十八年正月（1661年）}顺治帝去世，八岁玄烨即位为康熙帝，索尼、苏克萨哈、遏必隆、鳌拜辅政。地方层面应追问此事如何影响户口、迁徙、身份、劳役、土地或社区安全，而不是把行政分类当作群体自身的说明。核心来源为《清世祖实录》顺治十八年正月；《清史稿·圣祖本纪》；《康熙起居注》与内国史院遗诏档。编年记录定位朝廷所知的时间，不足以呈现城市、乡村、妇女、儿童和流动人口的全部经验。账本的研究方向是幼主继承、四辅政大臣与八旗政治。可核条件为顺治早逝使未成年皇帝继承，宗室与旗人高层需建立可承认的集体决策安排。应分层追问的后果为辅政集团掌握朝政，皇权成长、旗人权力平衡和南方战事资源相互交织。证据界限：原有置信注记：继位及辅政名单可靠；顺治死因和遗诏形成过程存在材料与后世叙事差异
\hypertarget{social:EQ-1661-COAST-EVACUATION}{}\paragraph{顺治十八年（1661年）}清廷下令福建广东沿海迁界，试图切断郑氏的粮食、人口和贸易联系。地方层面应追问此事如何影响户口、迁徙、身份、劳役、土地或社区安全，而不是把行政分类当作群体自身的说明。核心来源为《清世祖实录》顺治十八年迁界谕旨；《清史稿·食货志》；《福建通志》《广东通志》及沿海县志迁界条。编年记录定位朝廷所知的时间，不足以呈现城市、乡村、妇女、儿童和流动人口的全部经验。账本的研究方向是迁界政策、沿海社区与反郑战略。可核条件为郑氏水师能从沿岸获取补给和信息，清廷缺乏有效水师时选择强制迁徙制造隔离带。应分层追问的后果为渔盐、耕作和宗族网络受破坏；迁界成为海防政策的社会成本，也未完全阻绝海上往来。证据界限：原有置信注记：迁界命令可靠；实际迁撤距离、死亡与秘密贸易须以各县档案和方志核实
\hypertarget{social:EQ-1661-ZHENG-TAIWAN-LANDING}{}\paragraph{永历十五年／顺治十八年四月（1661年）}郑成功军登陆台湾并围攻荷兰东印度公司在热兰遮城的据点。地方层面应追问此事如何影响户口、迁徙、身份、劳役、土地或社区安全，而不是把行政分类当作群体自身的说明。核心来源为《清世祖实录》顺治十八年海防军报；《清史稿·郑成功传》；《热兰遮城日志》荷文1661年四月至十二月条。编年记录定位朝廷所知的时间，不足以呈现城市、乡村、妇女、儿童和流动人口的全部经验。账本的研究方向是郑氏转进台湾、荷兰殖民据点与原住民政治。可核条件为南京北伐失败后郑氏需安全粮食和军政基地，台湾荷兰据点又处海上航路和移民社会交汇处。应分层追问的后果为战争改变台湾政权框架与土地征发，东宁政权形成；原住民处境不能由任一方单述。证据界限：原有置信注记：登陆与围城可靠；兵员、原住民联盟与双方伤亡需区分中荷各自战报
\hypertarget{social:EQ-1662-ZEELANDIA}{}\paragraph{永历十六年／康熙元年正月（1662年）}热兰遮城守军投降，荷兰东印度公司在台湾西南的统治终结。地方层面应追问此事如何影响户口、迁徙、身份、劳役、土地或社区安全，而不是把行政分类当作群体自身的说明。核心来源为《清圣祖实录》康熙元年台湾相关军报；《清史稿·郑成功传》；《热兰遮城日志》荷文1662年正月及投降条约。编年记录定位朝廷所知的时间，不足以呈现城市、乡村、妇女、儿童和流动人口的全部经验。账本的研究方向是东宁建国、殖民交接与台湾多族群社会。可核条件为长期围城切断荷兰守军补给，郑氏控制周边村社与航道后取得谈判优势。应分层追问的后果为郑氏接收城堡、田地和贸易节点，台湾成为反清政权的财政与人口基地。证据界限：原有置信注记：投降文本和日期可靠；城外居民、原住民与被俘人员经历须从不同语言材料重建
\hypertarget{social:EQ-1662-YONGLI-EXECUTION}{}\paragraph{康熙元年四月（1662年）}吴三桂军将自缅甸押回的永历帝朱由榔处死，南明朝廷名义中心结束。地方层面应追问此事如何影响户口、迁徙、身份、劳役、土地或社区安全，而不是把行政分类当作群体自身的说明。核心来源为《清圣祖实录》康熙元年四月；《清史稿·永历帝本纪》；《琉璃宫史》缅文传统与《云南通志》。编年记录定位朝廷所知的时间，不足以呈现城市、乡村、妇女、儿童和流动人口的全部经验。账本的研究方向是永历末路、中缅边境交涉与南明合法性。可核条件为永历朝长期依赖流亡和边地补给，缅甸宫廷在清军压力下改变安置策略，吴三桂完成追捕。应分层追问的后果为南明皇帝名义不复存在，但西南残部、郑氏与地方记忆仍继续影响清初整合。证据界限：原有置信注记：被押回和处死主干可靠；缅甸移交过程、处死地点与责任叙事须比对中缅材料
\hypertarget{social:EQ-1662-ZHENG-DEATH}{}\paragraph{永历十六年／康熙元年五月（1662年）}郑成功在台湾去世，郑氏集团进入继承与军政重组阶段。地方层面应追问此事如何影响户口、迁徙、身份、劳役、土地或社区安全，而不是把行政分类当作群体自身的说明。核心来源为《清圣祖实录》康熙元年海防军报；《清史稿·郑成功传》；《巴达维亚城日记》1662年台湾消息与郑氏文书辑录。编年记录定位朝廷所知的时间，不足以呈现城市、乡村、妇女、儿童和流动人口的全部经验。账本的研究方向是郑成功之死、东宁继承与海上政权韧性。可核条件为热兰遮投降后郑氏尚未稳定整合台湾，郑成功骤死使官兵、宗族和海商网络重谈权力。应分层追问的后果为郑经最终掌权，东宁继续保有对抗清廷的海上基地，但内部资源争夺削弱集中行动。证据界限：原有置信注记：死亡月份大体可靠；死因、继承人之争与军中处置依赖郑氏后出材料和荷兰报告
\hypertarget{social:EQ-1663-XIAMEN-JINMEN}{}\paragraph{康熙二年（1663年）}清军联合部分郑氏降将攻取厦门、金门，郑氏主力退守台湾。地方层面应追问此事如何影响户口、迁徙、身份、劳役、土地或社区安全，而不是把行政分类当作群体自身的说明。核心来源为《清圣祖实录》康熙二年福建军报；《清史稿·施琅传》；《巴达维亚城日记》1663年中国沿海消息。编年记录定位朝廷所知的时间，不足以呈现城市、乡村、妇女、儿童和流动人口的全部经验。账本的研究方向是厦金失守、降将网络与清廷海防。可核条件为郑成功死后郑氏继承未稳，清军利用熟悉海战的降将和沿海据点进攻。应分层追问的后果为郑氏失去近大陆前进基地，台湾战略地位上升，清廷仍未取得制海能力。证据界限：原有置信注记：厦金易手可靠；荷兰参与、海战细节和郑氏损失应区分清方军报与荷兰记录
\hypertarget{social:EQ-1664-KUIZHOU}{}\paragraph{康熙三年（1664年）}清军攻破夔东据点，李来亨等抗清武装失败，长江上游的明遗民军事中心瓦解。地方层面应追问此事如何影响户口、迁徙、身份、劳役、土地或社区安全，而不是把行政分类当作群体自身的说明。核心来源为《清圣祖实录》康熙三年湖广四川军报；《清史稿·李来亨传》；《夔州府志》与南明遗民文集。编年记录定位朝廷所知的时间，不足以呈现城市、乡村、妇女、儿童和流动人口的全部经验。账本的研究方向是夔东抗清、山地据点与长江上游交通。可核条件为永历朝覆灭后夔东武装仍凭山地、峡江与地方网络抵抗，清廷需清理长江航运安全隐患。应分层追问的后果为清廷加强川东鄂西通道控制，残余反清力量更难与外部政权形成连线。证据界限：原有置信注记：夔东结局可靠；各部关系、山地据点与居民伤亡材料零散
\hypertarget{social:EQ-1665-REGENTS-CONSOLIDATION}{}\paragraph{康熙四年（1665年）}四辅政大臣以议政和旗务网络处理军政事务，鳌拜在政务中的影响扩大。地方层面应追问此事如何影响户口、迁徙、身份、劳役、土地或社区安全，而不是把行政分类当作群体自身的说明。核心来源为《清圣祖实录》康熙四年政务条；《清史稿·鳌拜传》；《康熙起居注》及内阁大库满汉题本。编年记录定位朝廷所知的时间，不足以呈现城市、乡村、妇女、儿童和流动人口的全部经验。账本的研究方向是辅政政治、议政机制与后设鳌拜叙事。可核条件为幼帝尚未亲裁，辅政大臣需在旗人贵族、部院官僚和南方军费之间协调。应分层追问的后果为中央决策呈现集体辅政与个人网络并存，为1669年权力重组提供背景。证据界限：原有置信注记：辅政架构可靠；权力高下多由后来的鳌拜案档与《清史稿》倒叙，须避免预设结局
\hypertarget{social:EQ-1667-SONI-DEATH}{}\paragraph{康熙六年（1667年）}首辅索尼去世，四辅政大臣的均衡被打破。地方层面应追问此事如何影响户口、迁徙、身份、劳役、土地或社区安全，而不是把行政分类当作群体自身的说明。核心来源为《清圣祖实录》康熙六年；《清史稿·索尼传》；《康熙起居注》康熙六年丧仪与议政记录。编年记录定位朝廷所知的时间，不足以呈现城市、乡村、妇女、儿童和流动人口的全部经验。账本的研究方向是索尼去世、辅政均衡与幼帝政治空间。可核条件为索尼居于辅政资历与旗人关系枢纽，其去世改变了原有制衡条件。应分层追问的后果为鳌拜及盟友影响更显著，康熙帝近侍和太皇太后方面开始寻找收回裁决权的机会。证据界限：原有置信注记：索尼死亡及官职变动可靠；其政治立场与遗言多经后出人物传记修饰
\hypertarget{social:EQ-1669-OBOI-ARREST}{}\paragraph{康熙八年五月（1669年）}康熙帝拘捕鳌拜并审理其党羽，结束辅政大臣实际主政的局面。地方层面应追问此事如何影响户口、迁徙、身份、劳役、土地或社区安全，而不是把行政分类当作群体自身的说明。核心来源为《清圣祖实录》康熙八年五月；《清史稿·鳌拜传》；《康熙起居注》与鳌拜案内阁档。编年记录定位朝廷所知的时间，不足以呈现城市、乡村、妇女、儿童和流动人口的全部经验。账本的研究方向是康熙亲政、旗人贵族与政治清算。可核条件为索尼死后辅政均衡失去支点，鳌拜与苏克萨哈案件激化皇帝近臣、太皇太后和旗人贵族冲突。应分层追问的后果为皇帝可直接统摄人事军政，但仍须通过议政、内阁及旗务结构行使权力。证据界限：原有置信注记：拘捕与处分可靠；宫中擒捕过程、罪状与少年皇帝个人谋略含胜利者叙事
\subsection{1670—1679：地方接收、军事移动与社会证据}
下列节点按账本顺序排列。每项政策、战役或制度变化，都应与其在家庭、劳动、道路和身份分类上的不均衡后果区分开来。
\hypertarget{social:EQ-1670-SACRED-EDICT}{}\paragraph{康熙九年（1670年）}康熙帝颁布十六条圣谕，规定由地方官和民众遵循的伦理、赋役与秩序训导框架。地方层面应追问此事如何影响户口、迁徙、身份、劳役、土地或社区安全，而不是把行政分类当作群体自身的说明。核心来源为《清圣祖实录》康熙九年十月；《清史稿·圣祖本纪》；《大清会典事例》圣谕条与各省地方志宣讲记录。编年记录定位朝廷所知的时间，不足以呈现城市、乡村、妇女、儿童和流动人口的全部经验。账本的研究方向是圣谕下乡、国家训导与地方中介。可核条件为清廷在征服后寻求以共同伦理语言补足武力和官僚的有限覆盖，地方官绅被要求传达。应分层追问的后果为圣谕成为后续地方教化与政治合法性资源，其接受和改写具有地区差异。证据界限：原有置信注记：颁布可靠；宣讲频率、识读范围和实际社会效果不能由文本证明
\hypertarget{social:EQ-1673-FEUDATORY-ABOLITION}{}\paragraph{康熙十二年三月（1673年）}康熙帝批准撤除吴三桂、尚可喜、耿继茂三藩的请求或安排，触发藩镇利益危机。地方层面应追问此事如何影响户口、迁徙、身份、劳役、土地或社区安全，而不是把行政分类当作群体自身的说明。核心来源为《清圣祖实录》康熙十二年三月；《清史稿·三藩传》；内阁大库康熙十二年撤藩题本与《平定三逆方略》。编年记录定位朝廷所知的时间，不足以呈现城市、乡村、妇女、儿童和流动人口的全部经验。账本的研究方向是撤藩决策、军费与中央—藩镇关系。可核条件为三藩长期占有军队、地盘和饷源，中央财政与皇权难以承受其永久化；尚可喜请撤使问题进入朝议。应分层追问的后果为撤藩将利益冲突公开化，吴三桂等选择武装反抗，清廷面临自入关以来最大内战。证据界限：原有置信注记：撤藩决策可靠；各藩请求先后、朝议立场和财政核算须以奏折及部议分开处理
\hypertarget{social:EQ-1673-WU-REVOLT}{}\paragraph{康熙十二年十一月（1673年）}吴三桂在云南起兵反清，三藩之乱正式爆发。地方层面应追问此事如何影响户口、迁徙、身份、劳役、土地或社区安全，而不是把行政分类当作群体自身的说明。核心来源为《清圣祖实录》康熙十二年十一月；《清史稿·吴三桂传》；《平定三逆方略》卷首及云南地方志。编年记录定位朝廷所知的时间，不足以呈现城市、乡村、妇女、儿童和流动人口的全部经验。账本的研究方向是吴三桂起兵、云南驻军与反清动员。可核条件为撤藩将削弱吴三桂的世袭军政和资源控制，云南驻军与西南道路使其拥有先发优势。应分层追问的后果为战火由云南扩展至华中、东南和西北，清廷需重建跨省饷运、指挥和地方合作。证据界限：原有置信注记：起兵及时间可靠；檄文中的族群、忠明和财政诉求是动员语言，不能直接等同参与者动机
\hypertarget{social:EQ-1674-GENG-JINGZHONG}{}\paragraph{康熙十三年（1674年）}靖南王耿精忠在福建反清，东南海陆战场与郑氏关系重新组合。地方层面应追问此事如何影响户口、迁徙、身份、劳役、土地或社区安全，而不是把行政分类当作群体自身的说明。核心来源为《清圣祖实录》康熙十三年福建军报；《清史稿·耿精忠传》；《平定三逆方略》与《福建通志》。编年记录定位朝廷所知的时间，不足以呈现城市、乡村、妇女、儿童和流动人口的全部经验。账本的研究方向是福建三藩、郑氏东宁与海陆联盟。可核条件为撤藩冲击福建藩镇利益，耿精忠既需郑氏海上支持又要维持陆上府县服从。应分层追问的后果为福建成为内战与海疆战争叠合战区，清廷行政和海防资源承受双重压力。证据界限：原有置信注记：耿精忠反叛可靠；其与郑氏盟约和实际协同行动须逐项辨析
\hypertarget{social:EQ-1674-WANG-FUCHEN}{}\paragraph{康熙十三年（1674年）}陕西提督王辅臣据平凉反清，西北军政和通往甘肃的交通受威胁。地方层面应追问此事如何影响户口、迁徙、身份、劳役、土地或社区安全，而不是把行政分类当作群体自身的说明。核心来源为《清圣祖实录》康熙十三年陕西军报；《清史稿·王辅臣传》；《平定三逆方略》与《平凉府志》。编年记录定位朝廷所知的时间，不足以呈现城市、乡村、妇女、儿童和流动人口的全部经验。账本的研究方向是平凉兵变、西北饷道与绿营军政。可核条件为撤藩战争使中央兵力南调，王辅臣对饷项、统属和地方资源的矛盾在西北爆发。应分层追问的后果为清廷面临关中至甘肃通道受阻，必须以招抚、围困和调兵稳定西北后方。证据界限：原有置信注记：平凉反叛可靠；王辅臣与吴三桂联络程度、兵数和民众态度不宜由檄文推定
\hypertarget{social:EQ-1674-ZHENG-JING-FUJIAN}{}\paragraph{康熙十三年（1674年）}郑经自台湾出兵福建，取得漳泉沿海若干据点并介入三藩战争。地方层面应追问此事如何影响户口、迁徙、身份、劳役、土地或社区安全，而不是把行政分类当作群体自身的说明。核心来源为《清圣祖实录》康熙十三年福建海防军报；《清史稿·郑成功传附郑经》；《台湾外纪》与《福建通志》。编年记录定位朝廷所知的时间，不足以呈现城市、乡村、妇女、儿童和流动人口的全部经验。账本的研究方向是郑经大陆行动、东宁财政与三藩内战。可核条件为清廷与耿精忠冲突削弱福建防线，郑氏试图以大陆税粮和港口减轻台湾资源压力。应分层追问的后果为郑氏短期扩大闽南影响，但陆上补给和盟友不稳使其难以转化为持久省域统治。证据界限：原有置信注记：出兵和若干据点易手可靠；郑耿关系、控制范围和征粮情况应以地方志与荷兰/商船信息校核
\hypertarget{social:EQ-1675-SHANG-ZHIXIN}{}\paragraph{康熙十四年（1675年）}平南王尚之信在广州控制军政并与反清势力周旋，两广局势动荡。地方层面应追问此事如何影响户口、迁徙、身份、劳役、土地或社区安全，而不是把行政分类当作群体自身的说明。核心来源为《清圣祖实录》康熙十四年广东军报；《清史稿·尚之信传》；《平定三逆方略》与《广东通志》。编年记录定位朝廷所知的时间，不足以呈现城市、乡村、妇女、儿童和流动人口的全部经验。账本的研究方向是两广藩镇、广州财政与政治摇摆。可核条件为撤藩冲击平南藩家族继承和军饷，三藩战争使尚之信可在清廷、吴军和地方势力间讨价还价。应分层追问的后果为两广不能成为清廷稳定后方，中央需同时处理招抚、军费和广东港口控制。证据界限：原有置信注记：尚之信的反复与广州控制大体可靠；是否完全反清须按不同时段命令、行动和归降分别记录
\hypertarget{social:EQ-1676-GENG-SURRENDER}{}\paragraph{康熙十五年（1676年）}耿精忠向清廷投降，福建藩镇叛乱的核心力量瓦解。地方层面应追问此事如何影响户口、迁徙、身份、劳役、土地或社区安全，而不是把行政分类当作群体自身的说明。核心来源为《清圣祖实录》康熙十五年；《清史稿·耿精忠传》；内阁大库康熙十五年福建招抚题本与《福建通志》。编年记录定位朝廷所知的时间，不足以呈现城市、乡村、妇女、儿童和流动人口的全部经验。账本的研究方向是招抚政策、福建战后军政与降军安置。可核条件为清廷陆路推进、海防压力和耿精忠内部离心共同作用，促成赦降条件。应分层追问的后果为福建大部重新纳入清方行政，但郑氏仍据台湾，降军安置成为新风险。证据界限：原有置信注记：投降与处分可靠；部众去向、郑氏影响和福建地方恢复须区别官方赦令与实际情况
\hypertarget{social:EQ-1676-WUCHANG}{}\paragraph{康熙十五年（1676年）}清军收复武昌等长江中游要地，吴军向湖南、广西方向收缩。地方层面应追问此事如何影响户口、迁徙、身份、劳役、土地或社区安全，而不是把行政分类当作群体自身的说明。核心来源为《清圣祖实录》康熙十五年湖广军报；《清史稿·赵良栋传》；《平定三逆方略》与《武昌府志》。编年记录定位朝廷所知的时间，不足以呈现城市、乡村、妇女、儿童和流动人口的全部经验。账本的研究方向是长江中游战局、漕运节点与清军反攻。可核条件为吴军北扩后补给线过长，清廷集结八旗、绿营和地方资源夺回江汉交通枢纽。应分层追问的后果为长江中游通道恢复有利于清军调度，吴三桂的战场重心被压回西南。证据界限：原有置信注记：武昌易手可靠；战役阶段、民众遭遇和双方控制线需按府县而非一城外推
\hypertarget{social:EQ-1677-WU-ZHOU}{}\paragraph{康熙十六年（1677年）}吴三桂在衡阳称帝，国号周，将撤藩战争转为公开争夺最高政治名义。地方层面应追问此事如何影响户口、迁徙、身份、劳役、土地或社区安全，而不是把行政分类当作群体自身的说明。核心来源为《清圣祖实录》康熙十六年湖南军报；《清史稿·吴三桂传》；《平定三逆方略》及吴周檄文辑录。编年记录定位朝廷所知的时间，不足以呈现城市、乡村、妇女、儿童和流动人口的全部经验。账本的研究方向是吴周建号、战争合法性与地方响应。可核条件为连年作战后吴三桂需要更高政治名义整合将领、赋役和盟友，其控制区又面临资源消耗。应分层追问的后果为清廷将战争表述为讨伐僭号政权；吴周建制化也增加维持官僚和军费的负担。证据界限：原有置信注记：称帝可靠；礼制、官号与各地响应范围主要见敌对军报，需与吴周文书互证
\hypertarget{social:EQ-1678-WU-DEATH}{}\paragraph{康熙十七年八月（1678年）}吴三桂去世，吴世璠继承吴周政权。地方层面应追问此事如何影响户口、迁徙、身份、劳役、土地或社区安全，而不是把行政分类当作群体自身的说明。核心来源为《清圣祖实录》康熙十七年八月；《清史稿·吴三桂传》；《平定三逆方略》与湖南地方志。编年记录定位朝廷所知的时间，不足以呈现城市、乡村、妇女、儿童和流动人口的全部经验。账本的研究方向是吴三桂之死、战争继承与吴周联盟。可核条件为长期作战与疾病、年龄压力使吴三桂政权失去创始者，继承人难以复制其将领和地方网络。应分层追问的后果为吴周指挥与合法性受冲击，清廷可逐步切断其湖南、贵州、云南联系。证据界限：原有置信注记：死亡与继承可靠；死因、遗命和军中稳定程度多来自后出叙事，应避免戏剧化细节
\hypertarget{social:EQ-1679-HUNAN-RECOVERY}{}\paragraph{康熙十八年（1679年）}清军在湖南持续推进并夺回若干府县，吴周军向贵州、广西收缩。地方层面应追问此事如何影响户口、迁徙、身份、劳役、土地或社区安全，而不是把行政分类当作群体自身的说明。核心来源为《清圣祖实录》康熙十八年湖南军报；《清史稿·赵良栋传》；《平定三逆方略》及《湖南通志》。编年记录定位朝廷所知的时间，不足以呈现城市、乡村、妇女、儿童和流动人口的全部经验。账本的研究方向是湖南收复、军粮与战争中的地方社会。可核条件为吴周失去长江中游节点后补给困难，清军利用水陆运输与分兵推进巩固新占地区。应分层追问的后果为吴周资源腹地缩小，清廷却须承担战区驻防、赈恤和恢复赋役的长期成本。证据界限：原有置信注记：清军推进主干可靠；具体县城归属和居民逃亡须由地方志、契约和军报逐县核查
\subsection{1680—1689：地方接收、军事移动与社会证据}
下列节点按账本顺序排列。每项政策、战役或制度变化，都应与其在家庭、劳动、道路和身份分类上的不均衡后果区分开来。
\hypertarget{social:EQ-1680-SICHUAN-RECOVERY}{}\paragraph{康熙十九年（1680年）}清军清理四川通往云南贵州的要道并恢复主要府城控制。地方层面应追问此事如何影响户口、迁徙、身份、劳役、土地或社区安全，而不是把行政分类当作群体自身的说明。核心来源为《清圣祖实录》康熙十九年四川军报；《清史稿·傅弘烈传》；《四川通志》与内阁大库四川粮饷题本。编年记录定位朝廷所知的时间，不足以呈现城市、乡村、妇女、儿童和流动人口的全部经验。账本的研究方向是四川战后重建、移民叙事与军需。可核条件为四川在多次战争中道路、仓储和行政受损，清廷需保障西南反攻后方补给。应分层追问的后果为清军取得向云南进军通道；招民垦荒与赋役恢复成为长期、地区不均的政策问题。证据界限：原有置信注记：主要军事结果可靠；四川战后人口、耕地和流亡统计口径极不一致
\hypertarget{social:EQ-1681-KUNMING}{}\paragraph{康熙二十年十月（1681年）}清军攻入昆明，吴世璠自尽，三藩之乱的吴周政权终结。地方层面应追问此事如何影响户口、迁徙、身份、劳役、土地或社区安全，而不是把行政分类当作群体自身的说明。核心来源为《清圣祖实录》康熙二十年十月；《清史稿·吴世璠传》；《平定三逆方略》与《云南通志》。编年记录定位朝廷所知的时间，不足以呈现城市、乡村、妇女、儿童和流动人口的全部经验。账本的研究方向是三藩之乱终局、西南占领与战后处置。可核条件为清廷经多年饷运、招抚和多路推进切断吴周资源，吴世璠缺少稳定援军与退路。应分层追问的后果为中央收回三藩地区名义军政权，但云南贵州驻军、财政和地方恢复仍需长期经营。证据界限：原有置信注记：昆明失守和吴世璠死亡可靠；围城损失、降附与后续清剿不得视为同日完成
\hypertarget{social:EQ-1681-ZHENG-JING-DEATH}{}\paragraph{康熙二十年（1681年）}郑经去世，东宁发生继承冲突，郑克塽最终掌权。地方层面应追问此事如何影响户口、迁徙、身份、劳役、土地或社区安全，而不是把行政分类当作群体自身的说明。核心来源为《清圣祖实录》康熙二十年台湾情报；《清史稿·郑成功传附郑经》；《台湾外纪》与荷兰商馆残存通信。编年记录定位朝廷所知的时间，不足以呈现城市、乡村、妇女、儿童和流动人口的全部经验。账本的研究方向是东宁继承危机、台湾军政与清廷决策窗口。可核条件为郑经长期在大陆作战并承受台湾财政压力，其死亡使宗族、将领和文官围绕继承重组。应分层追问的后果为东宁短期政治不稳，清廷得以评估以水师进攻台湾的风险和机会。证据界限：原有置信注记：郑经死期和继承结果可靠；宫廷政变过程主要依郑氏后出材料与清方情报
\hypertarget{social:EQ-1682-SHI-LANG}{}\paragraph{康熙二十一年（1682年）}清廷任命施琅统率福建水师，筹备进攻台湾。地方层面应追问此事如何影响户口、迁徙、身份、劳役、土地或社区安全，而不是把行政分类当作群体自身的说明。核心来源为《清圣祖实录》康熙二十一年福建海防条；《清史稿·施琅传》；《康熙朝满文朱批奏折全译》康熙二十一年台湾军务奏折。编年记录定位朝廷所知的时间，不足以呈现城市、乡村、妇女、儿童和流动人口的全部经验。账本的研究方向是施琅、福建水师与攻台决策。可核条件为三藩平定和东宁继承危机减少内陆牵制，厦金旧将对航道与郑氏军制的了解提供条件。应分层追问的后果为清廷从迁界防御转向跨海进攻，水师训练、粮械和季风仍决定行动可行性。证据界限：原有置信注记：任命和筹备可靠；施琅战略构想与朝廷内部争论须以奏折、朱批和后出传记区别处理
\hypertarget{social:EQ-1683-PENGHU}{}\paragraph{康熙二十二年六月（1683年）}施琅水师在澎湖击败郑氏舰队，取得通往台湾本岛的制海优势。地方层面应追问此事如何影响户口、迁徙、身份、劳役、土地或社区安全，而不是把行政分类当作群体自身的说明。核心来源为《清圣祖实录》康熙二十二年六月；《清史稿·施琅传》；施琅《飞报澎湖大捷疏》与台湾地方志。编年记录定位朝廷所知的时间，不足以呈现城市、乡村、妇女、儿童和流动人口的全部经验。账本的研究方向是澎湖海战、季风航路与水师后勤。可核条件为清军利用季风窗口、福建补给和施琅对郑氏舰队的熟悉发动集中攻击，郑氏难以同时守澎湖与本岛。应分层追问的后果为郑氏失去海上屏障，台湾方面转入投降谈判；澎湖伤亡和战术细节仍需多源检验。证据界限：原有置信注记：澎湖战役结果可靠；舰船、兵员和死伤数字多来自胜方战报，须以郑氏和海域材料约束
\hypertarget{social:EQ-1683-TAIWAN-SURRENDER}{}\paragraph{康熙二十二年八月（1683年）}郑克塽向清廷投降，东宁政权结束。地方层面应追问此事如何影响户口、迁徙、身份、劳役、土地或社区安全，而不是把行政分类当作群体自身的说明。核心来源为《清圣祖实录》康熙二十二年八月；《清史稿·郑成功传附郑克塽》；施琅奏折与《台湾府志》。编年记录定位朝廷所知的时间，不足以呈现城市、乡村、妇女、儿童和流动人口的全部经验。账本的研究方向是东宁投降、台湾接收与军民安置。可核条件为澎湖失败切断东宁海防与外援预期，继承危机使郑氏难以继续集中抵抗。应分层追问的后果为清廷取得台湾政治接收权，后续设治、迁界调整和土地社会重组并非投降即告完成。证据界限：原有置信注记：投降与清廷接收主干可靠；投降条件、官兵安置和原住民、移民社会变化须分层记录
\hypertarget{social:EQ-1684-TAIWAN-PREFECTURE}{}\paragraph{康熙二十三年（1684年）}清廷设置台湾府及诸县，隶属福建省，开始以府县和驻防制度经营台湾。地方层面应追问此事如何影响户口、迁徙、身份、劳役、土地或社区安全，而不是把行政分类当作群体自身的说明。核心来源为《清圣祖实录》康熙二十三年台湾设治条；《清史稿·地理志》；《台湾府志》与福建省档案相关奏议。编年记录定位朝廷所知的时间，不足以呈现城市、乡村、妇女、儿童和流动人口的全部经验。账本的研究方向是台湾设治、府县制度与边界化治理。可核条件为东宁投降后清廷须在弃守与设治间权衡海防、税源、移民和驻军成本。应分层追问的后果为台湾纳入福建行政体系，但沿海、平原、内山和原住民社群与官府关系呈现不均衡接触。证据界限：原有置信注记：设府设县可靠；机构设置不等于全岛有效治理，原住民地区和内山控制尤须保留
\hypertarget{social:EQ-1684-SEA-BAN-LIFT}{}\paragraph{康熙二十三年（1684年）}清廷调整迁界与海禁政策，允许沿海居民复界并逐步恢复民间海贸。地方层面应追问此事如何影响户口、迁徙、身份、劳役、土地或社区安全，而不是把行政分类当作群体自身的说明。核心来源为《清圣祖实录》康熙二十三年复界海禁条；《清史稿·食货志》；《福建通志》《广东通志》沿海复界条。编年记录定位朝廷所知的时间，不足以呈现城市、乡村、妇女、儿童和流动人口的全部经验。账本的研究方向是迁界终止、海贸恢复与沿海社会。可核条件为台湾投降降低清廷对郑氏补给网络的担忧，长期迁界的税源、民生和治安代价亦需处理。应分层追问的后果为沿海耕作、渔盐和航运开始恢复，但贸易许可、海防和地方执行塑造不同港口开放节奏。证据界限：原有置信注记：政策转向可靠；复界日期、港口开放和走私变化因省县而异，不能由中央谕旨推全国
\hypertarget{social:EQ-1685-ALBAZIN-FIRST}{}\paragraph{康熙二十四年（1685年）}清军围攻并迫使俄方撤离阿尔巴津，黑龙江边境进入直接谈判前的军事对峙。地方层面应追问此事如何影响户口、迁徙、身份、劳役、土地或社区安全，而不是把行政分类当作群体自身的说明。核心来源为《清圣祖实录》康熙二十四年黑龙江军报；《清史稿·邦交志》；俄国西伯利亚衙门文件与尼布楚谈判材料。编年记录定位朝廷所知的时间，不足以呈现城市、乡村、妇女、儿童和流动人口的全部经验。账本的研究方向是阿尔巴津、黑龙江防务与中俄边疆信息。可核条件为俄国哥萨克据点进入清廷所称黑龙江防务范围，双方地理认知、补给路线和外交权限不一致。应分层追问的后果为第一次围攻未解决边界问题，清俄转向再驻防、再对峙与谈判准备。证据界限：原有置信注记：围城与撤离可靠；守军人数、伤亡和领土主张分别来自清俄两方，不能以一方战报定量
\hypertarget{social:EQ-1686-ALBAZIN-SECOND}{}\paragraph{康熙二十五年（1686年）}俄方重返阿尔巴津后，清军再次围困，双方在疾病、补给与外交压力下僵持。地方层面应追问此事如何影响户口、迁徙、身份、劳役、土地或社区安全，而不是把行政分类当作群体自身的说明。核心来源为《清圣祖实录》康熙二十五年黑龙江军报；《清史稿·邦交志》；俄国阿尔巴津驻军报告及城址考古报告。编年记录定位朝廷所知的时间，不足以呈现城市、乡村、妇女、儿童和流动人口的全部经验。账本的研究方向是阿尔巴津二次围困、边境疾病与谈判杠杆。可核条件为第一次撤离未形成正式边界文本，俄方重建据点，清军以围困和东北后勤施压。应分层追问的后果为军事僵持使双方认识到谈判必要性，边界条约的语言与地图问题被推到前台。证据界限：原有置信注记：再围困可靠；兵力、疾病死亡与停战时点应并读满文、俄文与考古材料
\hypertarget{social:EQ-1688-GALDAN-KHALKHA}{}\paragraph{康熙二十七年（1688年）}噶尔丹进攻喀尔喀，喀尔喀诸部南徙并向清廷求援，北方草原秩序改变。地方层面应追问此事如何影响户口、迁徙、身份、劳役、土地或社区安全，而不是把行政分类当作群体自身的说明。核心来源为《清圣祖实录》康熙二十七年蒙古军报；《清史稿·噶尔丹传》；《蒙古源流》与俄罗斯使团材料。编年记录定位朝廷所知的时间，不足以呈现城市、乡村、妇女、儿童和流动人口的全部经验。账本的研究方向是准噶尔—喀尔喀战争、蒙古迁徙与清廷介入。可核条件为草原盟主竞争、牧地与贸易路线冲突促使噶尔丹东进，喀尔喀领主以清廷支援寻求生存空间。应分层追问的后果为清廷获得介入喀尔喀政治契机，但须承担难民安置、盟旗关系和对准噶尔战争成本。证据界限：原有置信注记：入侵与迁徙主干可靠；部众规模、附属关系和“归附”表述须并读蒙文与满汉文
\hypertarget{social:EQ-1689-NERCHINSK}{}\paragraph{康熙二十八年七月（1689年）}清俄在尼布楚签订条约，以满文、拉丁文和俄文文本处理边界与阿尔巴津问题。地方层面应追问此事如何影响户口、迁徙、身份、劳役、土地或社区安全，而不是把行政分类当作群体自身的说明。核心来源为《清圣祖实录》康熙二十八年七月；《清史稿·邦交志》；《尼布楚条约》满文、拉丁文、俄文文本及俄国谈判档。编年记录定位朝廷所知的时间，不足以呈现城市、乡村、妇女、儿童和流动人口的全部经验。账本的研究方向是尼布楚条约、多语外交与边界文本。可核条件为阿尔巴津围困和俄国东扩造成边境危机，耶稣会译员和双方使团使多语谈判成为可能。应分层追问的后果为阿尔巴津问题获得暂时制度化处理，边界解释和贸易联系仍需经后续使节、地图和地方执行确认。证据界限：原有置信注记：签约和多文本存在可靠；各文本术语、地名对应和界线含义不完全一致，不能投射为精确现代边界
\subsection{1690—1699：地方接收、军事移动与社会证据}
下列节点按账本顺序排列。每项政策、战役或制度变化，都应与其在家庭、劳动、道路和身份分类上的不均衡后果区分开来。
\hypertarget{social:EQ-1690-ULAN-BUTUNG}{}\paragraph{康熙二十九年（1690年）}清军在乌兰布通与噶尔丹军交战，阻止其向内蒙古腹地深入但未立即消灭其军队。地方层面应追问此事如何影响户口、迁徙、身份、劳役、土地或社区安全，而不是把行政分类当作群体自身的说明。核心来源为《清圣祖实录》康熙二十九年八月；《清史稿·噶尔丹传》；《亲征平定朔漠方略》及俄国商旅报告。编年记录定位朝廷所知的时间，不足以呈现城市、乡村、妇女、儿童和流动人口的全部经验。账本的研究方向是乌兰布通之战、草原补给与战报叙事。可核条件为喀尔喀危机后噶尔丹持续东进，清廷须防其联络内蒙古并威胁京师北方通道。应分层追问的后果为清军守住战略通道却未完成歼灭，战争转为更依赖补给、盟友和追击时机的长期行动。证据界限：原有置信注记：战役与战略结果大体可靠；胜负判断、兵力和伤亡在方略与准噶尔相关材料中不一
\hypertarget{social:EQ-1691-DOLONNOR}{}\paragraph{康熙三十年（1691年）}康熙帝在多伦诺尔会盟，安排喀尔喀诸部入盟旗体系并调度对噶尔丹资源。地方层面应追问此事如何影响户口、迁徙、身份、劳役、土地或社区安全，而不是把行政分类当作群体自身的说明。核心来源为《清圣祖实录》康熙三十年五月；《清史稿·藩部传》；理藩院满蒙文档案中的喀尔喀盟旗与赈济文书。编年记录定位朝廷所知的时间，不足以呈现城市、乡村、妇女、儿童和流动人口的全部经验。账本的研究方向是多伦会盟、盟旗制度与蒙古政治能动性。可核条件为噶尔丹战争导致喀尔喀领主寻求清廷军事支持，清廷以礼仪、封爵和补给将安全合作制度化。应分层追问的后果为清廷在蒙古影响扩大，同时承担牧地、赈济和跨部冲突调停责任，统属关系仍需谈判。证据界限：原有置信注记：会盟和册封可靠；“归附”不是单向完成，盟旗编制、牧地与领主自主性须以蒙文档案追踪
\hypertarget{social:EQ-1695-GALDAN-CAMPAIGN}{}\paragraph{康熙三十四年（1695年）}康熙帝筹划并发动对噶尔丹的第二阶段远征，集中军粮、驼马和多路部队。地方层面应追问此事如何影响户口、迁徙、身份、劳役、土地或社区安全，而不是把行政分类当作群体自身的说明。核心来源为《清圣祖实录》康熙三十四年军务条；《清史稿·圣祖本纪》；《康熙朝满文朱批奏折全译》噶尔丹军需奏折。编年记录定位朝廷所知的时间，不足以呈现城市、乡村、妇女、儿童和流动人口的全部经验。账本的研究方向是亲征后勤、草原战争与帝国动员。可核条件为乌兰布通后噶尔丹仍有机动空间，清廷需把盟旗支持、内地粮饷和季节性草原条件组织为追击。应分层追问的后果为清军以更大后勤投入缩小噶尔丹行动范围，也暴露远征对运输、疫病和环境的依赖。证据界限：原有置信注记：远征部署可靠；实际行军人数、牲畜损耗和地方征发须以军需档而非方略夸示数字判断
\hypertarget{social:EQ-1696-JAOMODO}{}\paragraph{康熙三十五年（1696年）}清军在昭莫多一带击败噶尔丹主力，噶尔丹向西遁走。地方层面应追问此事如何影响户口、迁徙、身份、劳役、土地或社区安全，而不是把行政分类当作群体自身的说明。核心来源为《清圣祖实录》康熙三十五年五月；《清史稿·噶尔丹传》；《亲征平定朔漠方略》与俄国边境报告。编年记录定位朝廷所知的时间，不足以呈现城市、乡村、妇女、儿童和流动人口的全部经验。账本的研究方向是昭莫多会战、满蒙协同与战果统计。可核条件为清军多路合围并获得喀尔喀和内蒙古盟友的情报、向导与补给，噶尔丹部众已受长期迁徙消耗。应分层追问的后果为噶尔丹军事核心受损，准噶尔对喀尔喀的威胁减弱，但战争未以一次会战完全结束。证据界限：原有置信注记：战役结局可靠；会战位置、伤亡与各部贡献应并读满文军报、蒙古材料和俄方消息
\hypertarget{social:EQ-1697-GALDAN-DEATH}{}\paragraph{康熙三十六年（1697年）}噶尔丹在科布多一带去世，其部众分散或向准噶尔核心地区退却。地方层面应追问此事如何影响户口、迁徙、身份、劳役、土地或社区安全，而不是把行政分类当作群体自身的说明。核心来源为《清圣祖实录》康熙三十六年；《清史稿·噶尔丹传》；俄国西伯利亚档案与蒙古口传整理。编年记录定位朝廷所知的时间，不足以呈现城市、乡村、妇女、儿童和流动人口的全部经验。账本的研究方向是噶尔丹之死、准噶尔政治与草原难民。可核条件为昭莫多败后补给、盟友和机动力枯竭，噶尔丹无法重建跨部联盟。应分层追问的后果为清廷宣布北方战事告捷，准噶尔政权仍在，蒙古牧地与人口流动后果持续存在。证据界限：原有置信注记：死亡事实大体可靠；死因、最后地点和部众归属说法各异，清廷凯旋叙事须与俄蒙材料互校
\hypertarget{social:EQ-1698-KHALKHA-ADMINISTRATION}{}\paragraph{康熙三十七年（1698年）}清廷在噶尔丹败亡后调整喀尔喀盟旗、台站与赈济安排，将战时同盟转为日常治理。地方层面应追问此事如何影响户口、迁徙、身份、劳役、土地或社区安全，而不是把行政分类当作群体自身的说明。核心来源为《清圣祖实录》康熙三十七年蒙古条；《清史稿·藩部传》；理藩院满蒙文档案中的喀尔喀盟旗与赈济文书。编年记录定位朝廷所知的时间，不足以呈现城市、乡村、妇女、儿童和流动人口的全部经验。账本的研究方向是战后蒙古治理、牧地与理藩院文书。可核条件为战争遗留难民、牲畜损失和权力真空，迫使清廷将军事保护转为行政、礼仪和物资安排。应分层追问的后果为盟旗体系获得更常态化清廷介入，但地方领主和牧民资源选择不能简化为被动纳入。证据界限：原有置信注记：制度调整可由谕令和理藩院文书确认；实际牧地使用、移徙和领主关系需长期追踪
\subsection{1700—1709：地方接收、军事移动与社会证据}
下列节点按账本顺序排列。每项政策、战役或制度变化，都应与其在家庭、劳动、道路和身份分类上的不均衡后果区分开来。
\hypertarget{social:EQ-1705-LHAZANG-TAKEOVER}{}\paragraph{康熙四十四年（1705年）}拉藏汗推翻第巴桑结嘉措的政权安排，西藏政教权力结构剧烈变化。地方层面应追问此事如何影响户口、迁徙、身份、劳役、土地或社区安全，而不是把行政分类当作群体自身的说明。核心来源为《清圣祖实录》康熙四十四年西藏条；《清史稿·西藏传》；《西藏王臣记》藏文与清廷驻藏来文。编年记录定位朝廷所知的时间，不足以呈现城市、乡村、妇女、儿童和流动人口的全部经验。账本的研究方向是拉藏汗政变、西藏政教关系与清廷信息链。可核条件为第巴与和硕特汗权长期紧张，达赖喇嘛继承和对外联络加深拉萨政治危机。应分层追问的后果为西藏内部权力重组为准噶尔介入和清廷后续军事行动创造条件，清方命令不等于拉萨实际控制。证据界限：原有置信注记：政变主干可靠；清廷知情、支持程度和地方反应须并读藏文、满文与旅行者材料
\hypertarget{social:EQ-1706-SIXTH-DALAI}{}\paragraph{康熙四十五年（1706年）}第六世达赖喇嘛仓央嘉措被押离拉萨，途中死亡或失踪，其结局与继承合法性成为争议。地方层面应追问此事如何影响户口、迁徙、身份、劳役、土地或社区安全，而不是把行政分类当作群体自身的说明。核心来源为《清圣祖实录》康熙四十五年西藏条；《清史稿·西藏传》；第巴传记、藏文史书与拉达克、蒙古相关记载。编年记录定位朝廷所知的时间，不足以呈现城市、乡村、妇女、儿童和流动人口的全部经验。账本的研究方向是第六世达赖喇嘛、转世认定与西藏政治叙事。可核条件为拉藏汗以宗教合法性重组拉萨权力，清廷册封语言与西藏各派、蒙古势力认定不一致。应分层追问的后果为达赖喇嘛地位争议削弱拉萨秩序，并为卡尔藏嘉措承认与清军入藏提供背景。证据界限：原有置信注记：被押离拉萨可靠；死亡地点、日期和是否在世的传统说法相互冲突，不能写成确定结论
\hypertarget{social:EQ-1708-KELZANG}{}\paragraph{康熙四十七年（1708年）}卡尔藏嘉措在青海塔尔寺附近被认定为第六世达赖喇嘛转世之一，围绕其身份形成跨区域支持网络。地方层面应追问此事如何影响户口、迁徙、身份、劳役、土地或社区安全，而不是把行政分类当作群体自身的说明。核心来源为《清圣祖实录》康熙四十七年青海西藏条；《清史稿·西藏传》；塔尔寺藏文档案与第七世达赖喇嘛传。编年记录定位朝廷所知的时间，不足以呈现城市、乡村、妇女、儿童和流动人口的全部经验。账本的研究方向是卡尔藏嘉措认定、青海寺院与转世政治。可核条件为第六世达赖喇嘛结局未决和拉藏汗所立人选缺乏广泛承认，青海寺院与蒙古关系提供替代合法性中心。应分层追问的后果为卡尔藏嘉措成为清廷、青海蒙古和西藏各派互动关键人物，宗教承认与军事保护相互牵连。证据界限：原有置信注记：出生与早期认定大体可靠；具体认定程序、时间和各方承认不可由后世转世谱系倒推
\subsection{1710—1719：地方接收、军事移动与社会证据}
下列节点按账本顺序排列。每项政策、战役或制度变化，都应与其在家庭、劳动、道路和身份分类上的不均衡后果区分开来。
\hypertarget{social:EQ-1711-TAX-FREEZE}{}\paragraph{康熙五十年（1711年）}康熙帝宣布以康熙五十年丁数为定额，后生人丁不再加征丁税。地方层面应追问此事如何影响户口、迁徙、身份、劳役、土地或社区安全，而不是把行政分类当作群体自身的说明。核心来源为《清圣祖实录》康熙五十年丁银谕；《清史稿·食货志》；户部档案与《大清会典事例》赋役条。编年记录定位朝廷所知的时间，不足以呈现城市、乡村、妇女、儿童和流动人口的全部经验。账本的研究方向是滋生人丁永不加赋、人口登记与地方财政。可核条件为长期战争后清廷寻求稳定赋役预期，人口登记与丁税征收的摩擦已显现。应分层追问的后果为政策改变丁额增长的制度预期，为雍正朝摊丁入亩等改革留下条件；实际减负须按省县账册验证。证据界限：原有置信注记：诏令可靠；各省执行时间、丁银与地丁并征口径和实际负担不能由诏令推定
\hypertarget{social:EQ-1712-TAX-EDICT}{}\paragraph{康熙五十一年（1712年）}清廷再次谕令固定丁额并要求地方遵行，围绕丁银征收与造册形成持续行政问题。地方层面应追问此事如何影响户口、迁徙、身份、劳役、土地或社区安全，而不是把行政分类当作群体自身的说明。核心来源为《清圣祖实录》康熙五十一年赋役谕；《清史稿·食货志》；户部题本与江苏、直隶等省赋役档。编年记录定位朝廷所知的时间，不足以呈现城市、乡村、妇女、儿童和流动人口的全部经验。账本的研究方向是丁税定额、命令传递与地方执行差异。可核条件为上一年诏令需经户部、督抚和州县造册落实，旧有缺额、摊派和财政压力造成执行摩擦。应分层追问的后果为“永不加赋”成为政策语言与后续改革依据，但登记人口、实际税负和地区差异不能被单一口号遮蔽。证据界限：原有置信注记：重复申谕可靠；它证明中央规范而非各地立即停止加派，需追踪地方账册与奏报
\hypertarget{social:EQ-1717-DZUNGAR-LHASA}{}\paragraph{康熙五十六年（1717年）}准噶尔军进入拉萨并杀拉藏汗，拉萨政局再次易手。地方层面应追问此事如何影响户口、迁徙、身份、劳役、土地或社区安全，而不是把行政分类当作群体自身的说明。核心来源为《清圣祖实录》康熙五十六年西藏军报；《清史稿·西藏传》；藏文《西藏王臣记》续编与准噶尔相关材料。编年记录定位朝廷所知的时间，不足以呈现城市、乡村、妇女、儿童和流动人口的全部经验。账本的研究方向是准噶尔入藏、拉萨权力真空与高原交通。可核条件为拉藏汗政权缺乏广泛稳定支持，准噶尔借宗教合法性与军事力量越过高原介入。应分层追问的后果为清廷面临保护卡尔藏嘉措、维护边疆通道和回应准噶尔扩张的复合压力，入藏决策加速。证据界限：原有置信注记：入拉萨和拉藏汗死亡可靠；暴力规模、民众态度与准噶尔治理需并读藏文、满文和旅行记录
\hypertarget{social:EQ-1718-QING-TIBET-EXPEDITION}{}\paragraph{康熙五十七年（1718年）}清军自四川方向入藏，在喀喇乌苏一带覆败，首批救援行动受挫。地方层面应追问此事如何影响户口、迁徙、身份、劳役、土地或社区安全，而不是把行政分类当作群体自身的说明。核心来源为《清圣祖实录》康熙五十七年川藏军报；《清史稿·西藏传》；《康熙朝满文朱批奏折全译》入藏军需奏折及藏文地方记载。编年记录定位朝廷所知的时间，不足以呈现城市、乡村、妇女、儿童和流动人口的全部经验。账本的研究方向是1718年入藏失败、高原后勤与信息不足。可核条件为清军低估高原道路、气候、补给与准噶尔防御，跨区域指挥链难以及时校正。应分层追问的后果为失败迫使清廷重组由青海、四川等方向的兵站与联盟，为1720年行动积累经验。证据界限：原有置信注记：远征失利可靠；路线、伤亡和藏地向导问题主要见清方军报，需与藏文地名和地方传统核对
\subsection{1720—1729：地方接收、军事移动与社会证据}
下列节点按账本顺序排列。每项政策、战役或制度变化，都应与其在家庭、劳动、道路和身份分类上的不均衡后果区分开来。
\hypertarget{social:EQ-1720-LHASA-RESTORATION}{}\paragraph{康熙五十九年（1720年）}清军及青海蒙古盟军进入拉萨，驱逐准噶尔力量并迎立卡尔藏嘉措为第七世达赖喇嘛。地方层面应追问此事如何影响户口、迁徙、身份、劳役、土地或社区安全，而不是把行政分类当作群体自身的说明。核心来源为《清圣祖实录》康熙五十九年西藏军报；《清史稿·西藏传》；满文军机前身档、藏文第七世达赖传与青海蒙古文书。编年记录定位朝廷所知的时间，不足以呈现城市、乡村、妇女、儿童和流动人口的全部经验。账本的研究方向是1720年入藏、清藏关系与多方军事联盟。可核条件为1718年失败后清廷改进补给并取得青海蒙古支持，准噶尔在拉萨的统治基础也不稳。应分层追问的后果为清廷在拉萨的军事政治存在增强，但西藏政教制度未被单一中央命令取代，后续治理仍待安排。证据界限：原有置信注记：进入拉萨和迎立主干可靠；清军、蒙古盟军与西藏地方行动者角色及伤亡不能由凯旋叙事抹平
\hypertarget{social:EQ-1721-ZHU-YIGUI}{}\paragraph{康熙六十年（1721年）}朱一贵在台湾发动反清起事，一度攻入府城，清廷随后派兵恢复控制。地方层面应追问此事如何影响户口、迁徙、身份、劳役、土地或社区安全，而不是把行政分类当作群体自身的说明。核心来源为《清圣祖实录》康熙六十年台湾军报；《清史稿·朱一贵传》；《台湾府志》、福建督抚奏折与契约文书。编年记录定位朝廷所知的时间，不足以呈现城市、乡村、妇女、儿童和流动人口的全部经验。账本的研究方向是朱一贵事件、台湾赋役与移民社会。可核条件为台湾设治后垦殖、租佃、族群关系与地方行政摩擦累积，官府征收和治安处置成为触发因素。应分层追问的后果为清廷暴露跨海调兵与府县治理薄弱环节，后续强化驻防和行政整顿；不能将起事简化为单一族群反应。证据界限：原有置信注记：起事、府城易手与镇压主干可靠；参与群体、赋税原因和伤亡数字受官府“逆匪”分类影响
\hypertarget{social:EQ-1722-KANGXI-DEATH}{}\paragraph{康熙六十一年十一月（1722年）}康熙帝去世，皇四子胤禛即位为雍正帝，早期清朝的长期统治转入新的继承阶段。地方层面应追问此事如何影响户口、迁徙、身份、劳役、土地或社区安全，而不是把行政分类当作群体自身的说明。核心来源为《清圣祖实录》康熙六十一年十一月；《清史稿·圣祖本纪、世宗本纪》；《雍正朝起居注册》与内阁大库遗诏档。编年记录定位朝廷所知的时间，不足以呈现城市、乡村、妇女、儿童和流动人口的全部经验。账本的研究方向是康熙继承、遗诏文本与雍正政治重组。可核条件为康熙晚年未公开稳定继承机制，宗室、旗人和官僚围绕皇位信息形成竞争与猜测。应分层追问的后果为雍正即位后将处理财政、旗务、边疆和宗室问题；传位细节不能以阴谋叙事代替可核文书。证据界限：原有置信注记：去世与即位日期可靠；遗诏、传位过程和储位争议受政治清算与后世传说影响，须以档案版本互校
